% Part: normal-modal-logic
% Chapter: frame-correspondence
% Section: equivalence-S5

\documentclass[../../../include/open-logic-section]{subfiles}

\begin{document}

\olfileid[hi]{nml}{frd}{es5}

\olsection{तुल्यता संबंध और \Log{S5}}

मोडल तर्कशास्त्र \Log{S5} को सभी सार्विक फ़्रेमों पर मान्य !!{formula}ों के
समुच्चय के रूप में अभिलक्षित किया जाता है; अर्थात् प्रत्येक जगत प्रत्येक
जगत से, स्वयं से भी, अभिगम्य है। ऐसी स्थिति में $\Box$ आवश्यकता और
$\Diamond$ संभावना के संगत हैं: यदि $!A$ \emph{प्रत्येक} जगत पर सत्य है तो
$\Box !A$ सत्य है, और यदि $!A$ \emph{किसी} जगत पर सत्य है तो $\Diamond !A$
सत्य है। यह निकलता है कि \Log{S5} को सभी स्वपरावर्ती, सममित और संक्रमणशील
फ़्रेमों, अर्थात् सभी \emph{तुल्यता संबंधों}, पर मान्य !!{formula}ों के रूप
में भी अभिलक्षित किया जा सकता है।

\begin{defn}
  $W$ पर कोई द्विपदी संबंध $R$ \emph{तुल्यता संबंध} है, तभी जब वह
  स्वपरावर्ती, सममित और संक्रमणशील हो। $W$ पर संबंध $R$ \emph{सार्विक} है,
  तभी जब प्रत्येक $u,v \in W$ के लिए $Ruv$।
\end{defn}

चूँकि \Ax{T}, \Ax{B} और \Ax{4} क्रमशः स्वपरावर्ती, सममित और संक्रमणशील
फ़्रेमों को अभिलक्षित करते हैं, इसलिए जिन फ़्रेमों में अभिगम्यता संबंध
तुल्यता संबंध है वे ठीक वही हैं जिनमें ये तीनों !!{formula} मान्य हैं। यह भी
निकलता है कि तुल्यता संबंधों को !!{formula}ों के अन्य संयोजनों से अभिलक्षित
किया जा सकता है, क्योंकि जिन प्रतिबंधों से हमने तुल्यता संबंध परिभाषित किए
हैं वे $R$ पर अन्य परिचित प्रतिबंधों के संयोजनों के तुल्य हैं।

\begin{prop}\ollabel{prop:equivalences}
  निम्नलिखित तुल्य हैं:
  \begin{enumerate}
  \item $R$ एक तुल्यता संबंध है;
  \item $R$ स्वपरावर्ती और यूक्लिडीय है;
  \item $R$ सीरियल, सममित और यूक्लिडीय है;
  \item $R$ सीरियल, सममित और संक्रमणशील है।
  \end{enumerate}
\end{prop}

\begin{proof}
  अभ्यास।
\end{proof}

\begin{prob}
  निम्नलिखित दिखाकर \olref[nml][frd][es5]{prop:equivalences} सिद्ध कीजिए:
  \begin{enumerate}
  \item यदि $R$ सममित और संक्रमणशील है, तो वह यूक्लिडीय है।
  \item यदि $R$ स्वपरावर्ती है, तो वह सीरियल है।
  \item यदि $R$ स्वपरावर्ती और यूक्लिडीय है, तो वह सममित है।
  \item यदि $R$ सममित और यूक्लिडीय है, तो वह संक्रमणशील है।
  \item यदि $R$ सीरियल, सममित और संक्रमणशील है, तो वह स्वपरावर्ती है।
  \end{enumerate}
  समझाइए कि प्रतिबंधों की तुल्यता के प्रमाण के लिए यह पर्याप्त क्यों है।
\end{prob}

\olref{prop:equivalences}, \olref[prf][prs]{prop:S5} का अर्थगत समकक्ष है,
क्योंकि यह उन फ़्रेमों के मोडल तर्कशास्त्र का तुल्य अभिलक्षण देता है जिन पर
$R$ तुल्यता संबंध है (परंपरागत रूप से \Log{S5} कहलाने वाला तर्कशास्त्र)।

सार्विक और तुल्यता संबंधों के बीच क्या संबंध है? यद्यपि प्रत्येक सार्विक
संबंध तुल्यता संबंध है, स्पष्टतः प्रत्येक तुल्यता संबंध सार्विक नहीं है।
किंतु सभी सार्विक संबंधों पर मान्य !!{formula} ठीक वही हैं जो सभी तुल्यता
संबंधों पर मान्य हैं।

\begin{prop}
  मान लें $R$ कोई तुल्यता संबंध है और प्रत्येक $w \in W$ के लिए $w$ का
  \emph{तुल्यता वर्ग} समुच्चय $[w] = \{w'\in W : Rww'\}$ से परिभाषित करें।
  तब:
  \begin{enumerate}
  \item $w \in [w]$;
  \item $R$ प्रत्येक तुल्यता वर्ग $[w]$ पर सार्विक है;
  \item तुल्यता वर्गों का संग्रह $W$ को परस्पर अपवर्जी और संयुक्त रूप से
    सर्वसमावेशी उपसमुच्चयों में विभाजित करता है।
  \end{enumerate}
\end{prop}

\begin{prop}\ollabel{prop:S5=univ}
  कोई !!{formula}~$!A$ उन सभी फ़्रेमों $\mModel{F} =\tuple{W,R}$ में मान्य है
  जहाँ $R$ तुल्यता संबंध है, तभी जब वह उन सभी फ़्रेमों
  $\mModel{F} =\tuple{W,R}$ में मान्य है जहाँ $R$ सार्विक है। अतः सार्विक
  फ़्रेमों का तर्कशास्त्र ठीक \Log{S5} है।
\end{prop}

\begin{proof}
  यह तत्काल जाँचा जा सकता है कि $W$ पर सार्विक संबंध~$R$ तुल्यता संबंध है।
  अतः यदि $!A$ उन सभी फ़्रेमों में मान्य है जहाँ $R$ तुल्यता संबंध है, तो वह
  सभी सार्विक फ़्रेमों में मान्य है। दूसरी दिशा के लिए प्रतिधनात्मक तर्क
  करें: मान लें $!B$ ऐसा !!{formula} है जो तुल्यता संबंध $R$ वाले फ़्रेम
  $\tuple{W,R}$ पर आधारित किसी प्रतिरूप
  $\mModel{M} = \tuple{W, R, V}$ के जगत $w$ पर विफल है। अतः
  $\mSat/{M}{!B}[w]$। प्रतिरूप $\mModel{M}' = \tuple{W',R',V'}$ इस प्रकार
  परिभाषित करें:
  \begin{enumerate}
  \item $W' = [w]$;
  \item $R'$, $W'$ पर सार्विक है;
  \item $V'(p) = V(p) \cap W'$।
  \end{enumerate}
  (अतः $\mModel{M}'$ के जगतों का समुच्चय $W'$,
  \olref{fig:partition} में छायांकित क्षेत्र से निरूपित है।) आसानी से देखा
  जा सकता है कि $R$ और $R'$, $W'$ पर सहमत हैं। तब !!{formula}ों पर आगमन से
  दिखाया जा सकता है कि प्रत्येक $!A$ और प्रत्येक $w' \in W'$ के लिए
  $\mSat{M'}{!A}[w']$ तभी जब $\mSat{M}{!A}[w']$ (यह सार्थक है क्योंकि
  $W' \subseteq W$)। विशेषतः $\mSat/{M'}{!B}[w]$, और $!B$ सार्विक फ़्रेम
  पर आधारित किसी प्रतिरूप में विफल है।
\end{proof}

\begin{figure}[t]
  \centering
  \begin{tikzpicture}[node distance=2cm, auto, thick]
    \clip [rounded corners] (0,0) -- (6,0) -- (6,4) -- (0,4) --  cycle;
    \begin{scope}
      \clip (0,0) -- (2,0)  .. controls (1.5,1.5) and (4.5,1.5)
      .. (4,4) -- (0,4) -- (0,0) -- cycle;
      \filldraw[fill=gray!40] (0,4) -- (0,2) .. controls (1.5,1.5)
      and (4.5,1.5) .. (4.5,0) -- (6,0) -- (6,4) -- (0,4) --  cycle;
    \end{scope}
    \draw [name path=line1] (2,0) .. controls (1.5,1.5) and (4.5,1.5) .. (4,4);
    \draw [name path=line2] (0,2)  .. controls (1.5,1.5) and (4.5,1.5) .. (4.5,0);
    \path [name intersections={of=line1 and line2}];
    \draw [rounded corners, use as bounding box, very thick] (0,0)
    -- (6,0) -- (6,4) -- (0,4) --  cycle; 
    \node at (2,3) {$[w]$} ;
    \node at (1,1) {$[u]$} ;
    \node at (3,0.75) {$[v]$} ;
    \node at (4.75,2) {$[z]$} ;
  \end{tikzpicture}
  \caption{$W$ का तुल्यता वर्गों में विभाजन।}
  \ollabel{fig:partition}
\end{figure}

\end{document}
